%%%%%%%%%%%%%%%%%%%%%%%%%%%%%%%%%%%%%%%%%%%%%%%%%%%%%%%%%%%
%\vspace{\sectionReduceTop}
\section{Results and Findings}
%\vspace{\sectionReduceBot}
%%%%%%%%%%%%%%%%%%%%%%%%%%%%%%%%%%%%%%%%%%%%%%%%%%%%%%%%%%%

%%%%%%%%%%%%%%%%%%%%%%%%%%%%%%%%%%%%%%%%%%%%%%%%%%%%%%%%%%%
\begin{table}
  \centering
  %\footnotesize
  \begin{tabular}{@{} p{0.5in} p{0.85in} cccccc@{}}
  %\begin{tabular}{@{} p{0.25in} p{0.6in} cccccc@{}}
      \toprule
      & & \multicolumn{2}{c}{Gibson} & \multicolumn{2}{c}{MP3D}\\
      \cmidrule(l){3-4} \cmidrule(l){5-6}
      Sensors & Baseline & SPL & Succ & SPL & Succ \\
      \midrule
      %\cmidrule{1-8}
      \multirow{4}{*}{\blind}
      & Random & $0.02$ & $0.03$ & $0.01$ & $0.01$ \\
      & Forward only & $0.00$ & $0.00$ & $0.00$ & $0.00$ \\
      & Goal follower & $0.23$ & $0.23$ & $0.12$ & $0.12$ \\
      & RL (PPO) & $0.42$ & $0.62$ & $0.25$ & $0.35$ \\
      \midrule
      %\cmidrule{2-8}
      \multirow{1}{*}{\rgb}
      & RL (PPO) & $0.46$ & $0.64$ & $0.30$ & $0.42$ \\
      \midrule
      %\cmidrule{2-8}
      \multirow{1}{*}{\depth}
      & RL (PPO) & $\textbf{0.79}$ & $\textbf{0.89}$ & $\textbf{0.54}$ & $\textbf{0.69}$ \\
      \midrule
      %\cmidrule{2-8}
      \multirow{2}{*}{\rgbd}
      & RL (PPO) & $0.70$ & $0.80$ & $0.42$ & $0.53$ \\
      & SLAM~\cite{mishkin2019benchmarking} & $0.51$ & $0.62$ & $0.39$ & $0.47$ \\
      \bottomrule
  \end{tabular}
  %\vspace{-2mm}
  \caption{Performance of baseline methods on the PointGoal task~\cite{Anderson2018-Evaluation} tested on the Gibson~\cite{Xia2018} and MP3D~\cite{Chang2017} test sets under multiple sensor configurations. RL models have been trained for $75$ million steps. We report average rate of episode success and SPL~\cite{Anderson2018-Evaluation}.
  }
  \label{tab:baselines}
  %\vspace{\captionReduceBot}
\end{table}

%%%%%%%%%%%%%%%%%%%%%%%%%%%%%%%%%%%%%%%%%%%%%%%%%%%%%%%%%%%


%Two key features of the Habitat stack are: i) generic dataset support and ii) high performance.
%These two features combined allow us to carry out experiments testing for generalization between datasets, and to explore training regimes with total experience far surpassing that of previous investigations.
We seek to answer two questions: 
i) how do learning-based agents compare to classic SLAM 
and hand-coded baselines as the amount of training experience increases and 
ii) how well do learned agents generalize 
across 3D datasets. 

It should be tacitly understood, but to be explicit -- 
`learning' and `SLAM' are broad families of techniques 
(and not a single method), 
are not necessarily mutually exclusive, and are 
not `settled' in their development. We compare representative instances of these families to gain insight into questions of 
scaling and generalization, and do not make any claims 
about intrinsic superiority of one or the other. 


\xhdr{Learning vs SLAM.}
To answer the first question we plot agent performance 
(\spl) on validation (\ie unseen) episodes over the 
course of training in \Cref{fig:classical_learning}
(top: Gibson, bottom: Matterport3D). 
%hows the results of this analysis. 
SLAM~\cite{mishkin2019benchmarking} 
does not require training 
and thus has a constant performance (0.59 on Gibson, 0.42 on 
Matterport3D). 
All RL (PPO) agents start out with far worse \spl, 
%In contrast, the RL (PPO) agents start out with far worse  performance.
%Over the course of training, 
but RL (PPO) \depth, in particular, improves dramatically 
and matches the classic baseline 
%once they have acquired 
at approximately $10$M frames 
(Gibson) or $30$M frames (Matterport3D)  
of experience, continuing to improve thereafter.  
Notice that if we terminated the experiment at $5$M frames 
as in \cite{mishkin2019benchmarking} we would 
also conclude that SLAM~\cite{mishkin2019benchmarking} dominates. 
Interestingly, 
%agents using only color vision (\rgb) as sensory input 
\rgb agents do not
significantly outperform \blind agents; we hypothesize 
%the agents using only the goal position input (\blind).
because both 
are equipped with GPS sensors. Indeed, qualitative 
results (\Cref{fig:top_down_nav} and video in supplement) 
suggest that \blind agents `hug' walls and 
implement `wall following' heuristics.
In contrast, \rgb sensors provide a high-dimensional complex 
signal that may be prone to overfitting to train environments 
due to the variety across scenes (even within the same dataset). 
%This might indicate that color-only agents are unable 
%to handle the visual complexity and variety of the environments in our datasets and thus do not generalize well to the unseen test environments (even from the same dataset).
We also notice in \Cref{fig:classical_learning} that \emph{all methods} perform better on Gibson than Matterport3D. 
This is consistent with our previous analysis that Gibson 
contains smaller scenes and shorter episodes. 
%scenes 
%are physically smaller and the episodes are shorter. 



Next, for each agent and dataset, we select the best-performing checkpoint on validation and report results on test in \Cref{tab:baselines}. 
%for all datasets. 
%A full summary of the best-performing agents evaluated on the test set is shown in \Cref{tab:baselines} for all datasets. (The agents were selected via SPL on the validation set.)
We observe that uniformly across the datasets, 
RL (PPO) \depth performs best, 
outperforming RL (PPO) \rgbd (by 0.09-0.16 \spl), 
SLAM (by 0.15-0.28 \spl), and \rgb (by 0.13-0.33 \spl) in that order (see the supplement for additional experiments involving noisy depth).
We believe \depth performs better than \rgbd because i) the PointGoal navigation task requires reasoning only about free space and depth provides relevant information directly, and ii) \rgb has significantly more entropy (different houses look very different), thus it is easier to overfit when using \rgb.
We ran our experiments with 5 random seeds per run, to confirm that these differences are statistically significant.
The differences are about an order of magnitude larger than the standard deviation of average SPL for all cases (\eg on the Gibson dataset errors are, \depth: $\pm 0.015$, \rgb: $\pm 0.055$, \rgbd: $\pm 0.028$, \blind: $\pm 0.005$).
%as measured by all metrics, on all datasets.
%\rgbd variant of the same agent and SLAM rank second or third, while the \rgb agent ranks fourth.
Random and forward-only agents have very low performance, while the hand-coded goal follower and \blind baseline see modest performance.%\footnotemark[4]
See the supplement for additional analysis of trained agent behavior.
%(approaching the performance of the \rgb in the latter case). 



%%%%%%%%%%%%%%%%%%%%%%%%%%%%%%%%%%%%%%%%%%%%%%%%%%%%%%%%%%%
\begin{figure}
    %\vspace{-10pt}
    \centering
    \includegraphics[width=\linewidth]{fig/top_down_nav.pdf}
    %\vspace{-20pt}
    \caption{Navigation examples for different sensory configurations of the RL (PPO) agent, visualizing trials from the Gibson and MP3D val sets. A \textbf{blue dot} and \textbf{red dot} indicate the starting and goal positions, and the \textbf{blue arrow} indicates final agent position. The \textbf{blue-green-red line} is the agent's trajectory. Color shifts from blue to red as the maximum number of agent steps is approached.
    See the supplemental materials for more example trajectories.
    }
    \label{fig:top_down_nav}
    %\vspace{\captionReduceBot}
    %\vspace{-5pt}
\end{figure}
%%%%%%%%%%%%%%%%%%%%%%%%%%%%%%%%%%%%%%%%%%%%%%%%%%%%%%%%%%%

In \Cref{fig:top_down_nav} we plot example trajectories for 
the RL (PPO) agents, to qualitatively contrast their 
behavior in the same episode. 
Consistent with the aggregate statistics, we observe that \blind collides with obstacles and follows walls, while \depth is 
the most efficient.
See the supplement and the video for more example trajectories.

%%%%%%%%%%%%%%%%%%%%%%%%%%%%%%%%%%%%%%%%%%%%%%%%%%%%%%%%%%%
\begin{figure}
    \centering
    \includegraphics[width=\columnwidth]{fig/generalization.pdf}
    %\vspace{-20pt}
    \caption{Generalization of agents between datasets. We report average SPL for a model trained on the source dataset in each row, as evaluated on test episodes for the target dataset in each column.}
    \label{fig:generalization}
    %\vspace{\captionReduceBot}
\end{figure}
%%%%%%%%%%%%%%%%%%%%%%%%%%%%%%%%%%%%%%%%%%%%%%%%%%%%%%%%%%%

\xhdr{Generalization across datasets.}
Our findings so far are that RL (PPO) agents significantly outperform SLAM~\cite{mishkin2019benchmarking}. This prompts our second question
-- are these findings dataset specific or do learned agents generalize across datasets?
%This is important because SLAM is (in principle) dataset agnostic. 
We report exhaustive comparisons in \Cref{fig:generalization} -- 
specifically, average \spl for all combinations of 
$\{$train, test$\} \times \{$Matterport3D, Gibson$\}$ for 
all agents $\{$\blind, \rgb, \rgbd, \depth$\}$. 
Rows indicate (agent, train set) pair, columns indicate test set. 
%in \Cref{fig:generalization}. 
%A summary of the generalization performance between datasets is provided in \Cref{fig:generalization} which reports average SPL on all pairs of datasets for train and test set.
We find a number of interesting trends. 
First, nearly all agents suffer a drop in performance 
when trained on one dataset and tested on another, 
\eg \rgbd Gibson$\rightarrow$Gibson $0.70$ vs \rgbd 
Gibson$\rightarrow$Matterport3D $0.53$ (drop of 0.17). 
\rgb and \rgbd agents suffer a 
significant performance degradation, 
while the \blind agent is least affected (as we would expect).

%We see a similar trend for agents tested on MP3D, though the performance degradation is much smaller.

Second, we find a potentially counter-intuitive trend -- 
agents trained on Gibson consistently outperform their 
counterparts trained on Matterport3D, \emph{even 
when evaluated on Matterport3D}.
We believe the reason is the previously noted observation that 
Gibson scenes are smaller and episodes are shorter 
(lower \gdsp) than Matterport3D. 
%A property of the Gibson dataset that might explain this observation is that it has a lower average geodesic distance (\gdsp) than MP3D (see dataset statistics in supplement).
Gibson agents are trained on `easier' episodes and 
encounter positive reward more easily during random exploration, 
thus bootstrapping learning. Consequently, for a fixed 
computation budget Gibson agents are stronger universally (not 
just on Gibson). 
This finding suggests that visual navigation agents could 
benefit from curriculum learning. 

These insights are enabled by the engineering of Habitat, which made these experiments as simple as a change in the evaluation dataset name.